\section{Ontologie HMT--MD des particules}
\Needspace{7\baselineskip}
\subsection{État généalogique matériel : extension, parcours, enregistrement, signature et représentation}
Un état HMT complet n'est pas un simple point. C'est le tuple

\[
 \mathfrak x=
 (\Gamma,\mathcal L_\Gamma,\sigma_\Gamma,\eta_\Gamma,
 \mathcal R_\Gamma,\mathcal B_\Gamma,\mathcal C_\Gamma),
\]
où \(\Gamma\) est un parcours survivant, \(\mathcal L_\Gamma\) son enregistrement généalogique, \(\sigma_\Gamma\) sa signature entière, \(\eta_\Gamma\) son raffinement harmonique, \(\mathcal R_\Gamma\) sa représentation, \(\mathcal B_\Gamma\) sa mémoire de retour et \(\mathcal C_\Gamma\) sa donnée de composition ou de frontière. Deux états ayant le même mot visible peuvent être distincts s'ils diffèrent par le préfixe, le bloc, la mémoire, la retenue, le feuillet ou l'ombre de Witt.

\Needspace{7\baselineskip}
\subsection{Particule comme section locale persistante de l'état généalogique dimensionnel}
Une particule est une classe d'équivalence de sections persistantes satisfaisant quatre conditions :

\begin{enumerate}
\item prolongement compatible à travers les cylindres du TPK ;
\item retour de phase avec mémoire d'état ;
\item représentation physique irréductible ou bloc invariant minimal ;
\item stabilité spectrale sous le couplage qui l'observe.
\end{enumerate}
Une section persistante est une famille compatible

\[
 s=(s_K)_{K\ge K_0},\qquad
 p_{L,K}(s_L)=s_K\quad(L\ge K\ge K_0).
\]
Deux sections \(s\) et \(s'\) représentent la même particule lorsqu'il existe un niveau \(K_1\) et une famille compatible de symétries internes \(g_K\in G_K\) tels que

\[
 s'_K=g_Ks_K\quad(K\ge K_1),
 \qquad
 p_{L,K}g_L=g_Kp_{L,K},
\]
et que les caractères observables de l'enregistrement généalogique coïncident. Cette relation identifie les descriptions qui diffèrent par la jauge interne, mais n'identifie pas les parcours de mémoire, retenue, feuillet, bloc ou holonomie distincts. La particule est la classe

\[
 [s]_{\rm pers}
 =\{s':s'\sim s\},
\]
et non un mot visible isolé ou une cellule choisie après connaissance de sa masse.

Si \(\widetilde\Gamma_K\) est l'extension au niveau \(K\), la persistance s'exprime par

\[
 p_{K+9,K}(\widetilde\Gamma_{K+9})
 =\widetilde\Gamma_K,\qquad
 g(K+9)=g(K),
\]
sans exiger

\[
 \mathcal B_{K+9}=\mathcal B_K.
\]
La particule n'est pas la phase qui revient, mais l'identité généalogique conservée à travers le changement d'état. Sa masse mesure la rigidité spectrale de cette persistance lorsqu'elle se couple à une réalisation matérielle.

\Needspace{7\baselineskip}
\subsection{Antiparticule comme image conjuguée d'une section persistante sous inversion de feuillet}
La conjugaison dodécaphasique est

\[
 C_M(\mathbf t)=-\operatorname{rev}(\mathbf t).
\]
La paire particule--antiparticule est formée de deux orientations d'une même bande généalogique. Si \(\chi\) désigne l'orientation de feuillet,

\[
 (\mathbf t,\chi)\sim(C_M\mathbf t,-\chi).
\]
Pour une énergie de bande

\[
 E_{\rm band}(\mathbf t,\chi)
 =E_+(\mathbf t)+\chi E_-(\mathbf t),
\]
avec

\[
 E_\pm(\mathbf t)
 =\frac12\bigl(E_0(\mathbf t)\pm E_0(C_M\mathbf t)\bigr),
\]
on obtient

\[
 E_{\rm band}(C_M\mathbf t,-\chi)
 =E_{\rm band}(\mathbf t,\chi).
\]
L'égalité des masses est donc une conséquence de la parité de l'opérateur sous conjugaison, tandis que charge et orientation changent. L'égalité des largeurs exige en outre que l'opérateur de désintégration soit conjugué par \(C_M\).

\Needspace{7\baselineskip}
\subsection{Masse comme caractère spectral de persistance de \(1\) parcours matériel}
La masse est la mesure spectrale, dans la droite \(\mathcal L_M\), de l'opérateur résultant de la compression du caractère, des tours, du lecteur bidirectionnel et de l'horloge sur le sous-espace persistant sélectionné par le couplage :

\[
 \widehat M_{P,a}
 =P_{P,a}\,
 \mathbf m_\star^{\rm ret}
 R_{\rm act}^{\widehat K_{P,a}/2}
 \widehat{\mathcal A}_{P,a}^{(0)}
 R_{\rm act}^{\widehat K_{P,a}/2}
 P_{P,a}.
\]
Le résultat général est

\[
 \mu_{P,a}^{\rm mass}(B)
 =\operatorname{Tr}\!\left(
 \rho_{P,a}E_{\widehat M_{P,a}}(B)\right).
\]
Il existe une masse scalaire \(m_P\) si et seulement si

\[
 P_{\operatorname{supp}\rho}\widehat M_{P,a}
 P_{\operatorname{supp}\rho}
 =m_P P_{\operatorname{supp}\rho},
\]
de manière équivalente, si la variance de \(\widehat M_{P,a}\) dans l'état préparé est nulle. Cette définition inclut l'électron de rang un, les multiplets imposés par Schur et les pôles isolés des résonances, sans les confondre.

\Needspace{7\baselineskip}
\subsection{Charge comme caractère de jauge transporté par la section matérielle}
La charge est le caractère orienté de la couche impaire de l'enregistrement généalogique sous conjugaison. Soit \(\sigma_{\rm odd}\) la signature impaire. Alors

\[
 Q(C_M\Gamma)=-Q(\Gamma),\qquad
 Q(\Gamma)=\mathfrak q\!\left(\sigma_{\rm odd}(\Gamma)\right),
\]
où \(\mathfrak q\) est l'homomorphisme de la signature impaire vers le réseau de charge de la réalisation. La signature massique paire et la signature de charge impaire sont des objets distincts : un mot peut avoir une somme linéaire nulle tout en portant une signature massique non nulle.

\Needspace{7\baselineskip}
\subsection{Spin comme représentation de l'action des rotations sur la fibre matérielle}
Le spin est la classe d'holonomie avec laquelle une section persistante répond à une rotation complète du repère d'orientation. La racine interne

\[
 J^2=-I
\]
sépare les deux clôtures :

\[
 U(2\pi)=
 \begin{cases}
 +I,&\text{secteur tensoriel},\\
 -I,&\text{secteur spinoriel},
 \end{cases}
 \qquad
 U(4\pi)=I.
\]
Un boson de spin entier appartient à une représentation qui se clôt à \(2\pi\) ; un fermion de spin demi-entier nécessite le double tour \(4\pi\). La valeur du spin ne se déduit pas d'un dénombrement cardinal isolé : c'est le poids de la représentation du groupe d'holonomie sur le sous-espace persistant. En particulier :

\[
 s=0:\text{ section scalaire},\qquad
 s=\tfrac12:\text{ bande spinorielle},
\]
\[
 s=1:\text{ section vectorielle transverse},\qquad
 s=2:\text{ composante tensorielle d'une réalisation métrique}.
\]
La dernière ligne définit un type de représentation ; elle n'impose pas l'existence d'une particule élémentaire de spin deux.

La chiralité et l'hélicité sont des données ultérieures et distinctes. Une involution orientée \(\Gamma_{\rm ch}^2=I\) définit

\[
 P_L=\frac{I-\Gamma_{\rm ch}}2,
 \qquad
 P_R=\frac{I+\Gamma_{\rm ch}}2,
\]
tandis que l'hélicité est le signe spectral du générateur de rotation suivant la direction de propagation. La première classe les feuillets de la représentation ; la seconde dépend de l'état cinématique. Cette séparation permet à un neutrino d'être fermionique par sa clôture \(4\pi\), même lorsque son couplage faible sélectionne une seule chiralité.

\Needspace{7\baselineskip}
\subsection{Statistique quantique induite par la symétrie d'échange des sections}
La statistique est la règle de complétion de l'espace des parcours identiques :

\[
 \mathcal F_-(\mathcal H)
 =\bigoplus_{n\ge0}\wedge^n\mathcal H,\qquad
 \mathcal F_+(\mathcal H)
 =\bigoplus_{n\ge0}\operatorname{Sym}^n\mathcal H.
\]
Dans la complétion extérieure, les relations CAR

\[
 \{a_i,a_j^\dagger\}=\delta_{ij},\qquad
 \{a_i,a_j\}=0
\]
impliquent

\[
 n_i^2=n_i
\]
et, par conséquent, l'exclusion de Pauli et la distribution de Fermi--Dirac. Dans la complétion symétrique, les relations CCR permettent une occupation multiple et conduisent à Bose--Einstein. La statistique ne s'infère pas de la masse : elle procède de la façon dont le parcours se compose sous échange.

\Needspace{7\baselineskip}
\subsection{Couleur comme représentation résiduelle de \(SU(3)\) sur les parcours de quarks}
La couleur est l'orbite résiduelle triadique d'une section de quark dans la fibre de charge et de saveur. Les trente-six positions colorées se décomposent comme

\[
 12_R+12_G+12_B.
\]
L'objet physique n'est pas l'une de ces couleurs, mais le triplet et son action :

\[
 \mathcal H_q^{\rm color}
 =\mathcal H_R\oplus\mathcal H_G\oplus\mathcal H_B.
\]
Si l'opérateur de masse commute avec l'action irréductible de couleur,

\[
 U_g\widehat M_qU_g^*=\widehat M_q,
\]
le lemme de Schur impose une masse commune dans le triplet. Choisir une cellule rouge, verte ou bleue comme représentante détruirait la covariance.

\Needspace{7\baselineskip}
\subsection{Saveur comme classe de parcours au sein d'une génération matérielle}
La saveur est une classe de couplage non équivalente au sein d'une même représentation de jauge, caractérisée par charge, génération, feuillet, signature fine, mémoire et projecteur d'interaction. Deux saveurs peuvent partager une fibre grossière et différer par l'entrelaceur qui les couple au lecteur.

Formellement, si \(F\) est la fibre commune, les saveurs \(f_i\) sont des représentations \(\pi_i\) pour lesquelles

\[
 \operatorname{Hom}_{G_F}(V_{f_i},\ell^2(F))\ne0
\]
et dont les projecteurs \(P_{f_i}\) ne sont pas conjugués par une symétrie de jauge qui demeure physique. La saveur n'est pas la profondeur de l'atlas et ne se réduit pas à un nom conventionnel.

\Needspace{7\baselineskip}
\subsection{Génération matérielle comme strate récurrente typée par profondeur de retour, signature fine et spectre du lecteur}
Une génération est une strate récurrente de la même grammaire de charge et de représentation, distinguée par profondeur de retour, signature fine et spectre du lecteur. L'étiquette de génération n'est physique que lorsqu'il existe un projecteur qui entrelace cette strate avec le canal de couplage. Une profondeur \(G_j\) ne doit donc pas être identifiée automatiquement à l'électron, au muon, au tau ou au neutrino stérile.

\Needspace{7\baselineskip}
\subsection{Champ de jauge comme connexion sur les fibres de charge}
Un champ de jauge est la réalisation d'une redondance de repère sur les fibres. Ses bosons élémentaires appartiennent à la direction nulle de l'opérateur de masse tant que la symétrie demeure exacte. Une masse de jauge ne s'obtient pas en approchant zéro par un caractère positif, mais par une brisure, un couplage longitudinal ou une réalisation de pôle déclarée.

\Needspace{7\baselineskip}
\subsection{Particule composée comme section persistante du produit de parcours liés}
Une particule composée est une section persistante de l'opérateur de composition, et non la somme des masses ou des signatures de ses constituants. Si \(\Gamma_1,\ldots,\Gamma_n\) sont les parcours constituants, on construit d'abord

\[
 \mathfrak C_{12}
 (\Gamma_1,\ldots,\Gamma_n;\mathcal T,\partial,\mathrm{carry})
 =\mathcal L_{\rm comp},
\]
où \(\mathcal T\) est l'arbre temporel et \(\partial\) la frontière de liaison. Ce n'est qu'ensuite qu'agit le lecteur

\[
 \mathcal L_{\rm comp}\longrightarrow\sigma_{\rm comp}
 \longrightarrow\chi_{\rm comp}
 \longrightarrow\widehat M_{\rm comp}.
\]
L'énergie de liaison réside dans le cocycle de composition, la mémoire de frontière et la modification du spectre ; elle n'est pas introduite comme correction externe.

\Needspace{7\baselineskip}
\subsection{Résonance comme section transitoire et largeur comme taux de désintégration}
Une résonance est un mode persistant non unitaire de l'opérateur couplé. Si

\[
 \lambda=re^{-i\theta},\qquad 0<r<1,
\]
alors

\[
 E=\frac{\hbar_{\rm int}}{t_0}
 \bigl(\theta-i\log r\bigr)
 =M-\frac{i}{2}\Gamma.
\]
La masse \(M\) et la largeur \(\Gamma\) sont les parties réelle et imaginaire d'un pôle, mais exigent des lecteurs distincts. La largeur peut aussi s'obtenir à partir d'un semi-groupe comprimé :

\[
 T_X=Q_XUQ_X,\qquad
 \lambda_X=-\frac1{t_0}\log\rho(T_X),\qquad
 \Gamma_X=\hbar_{\rm int}\lambda_X.
\]
On ne déduit pas une largeur du seul fait qu'un parcours possède un horizon fini.

\Needspace{7\baselineskip}
\subsection{Section virtuelle comme extension intermédiaire non asymptotique de l'atlas}
Une section est virtuelle lorsque son prolongement se termine de manière non triviale :

\[
 \operatorname{Ext}_{L}(s)\ne\varnothing,\qquad
 \operatorname{Ext}_{L+1}(s)=\varnothing,\qquad
 \mathcal L(s)=L<\infty.
\]
Son contenu primaire est l'horizon \(L\), l'obstruction terminale, l'enregistrement généalogique fini et l'amplitude de propagation. Elle n'est ni une particule de masse imaginaire ni un état documentaire non encore classé.

\Needspace{7\baselineskip}
